% chapters/08_uebergabe.tex

\chapter{Übergabe}
\label{chap:uebergabe}
\label{chap:fazit}
\label{chap:weiterfuehrend}
\label{sec:ausblick}
\label{sec:konsequenz}

\section{Fazit}
\label{sec:fazit}

Die drei Ziele aus \cref{sec:aufgabe} sind unterschiedlich weit erreicht.
(1) Das Fine-Tuning läuft reproduzierbar lokal und auf KISSKI; auf echten Bildern folgt das
Modell den Demonstrationen, N1.6 und N1.7 gleichauf (\cref{sec:erg-training}).
(2) Die Bewertung ohne Roboter steht, und sie erklärt das Scheitern in der Simulation mit
zwei gemessenen Ursachen: Würfellage und visueller Domain-Gap. Das Co-Training mit
gerenderten Bildern verkleinert den Gap und bringt erste Griffe (2 von 60), aber kein
Stapeln.
(3) Schnittstellen und Anleitung sind in \cref{chap:schnittstellen,chap:anleitung} und im
Repository dokumentiert, Versionen und Reproduzierbarkeit in \cref{anh:software}. Die Aufgabe selbst, drei Würfel zu stapeln, ist nicht gelöst. Ob das
Modell auf dem echten Roboter greift, ist offen; dieser Test steht an erster Stelle der
folgenden Liste.

\section{Offene Punkte, in dieser Reihenfolge}
\label{sec:offene-punkte}

Die laufend gepflegte Liste steht in \doc{next-steps.md}, Abschnitt~2.

\begin{enumerate}
  \item \textbf{Robotertest.} Modell C (\texttt{…-cotrain-v3-cube4-ckpt44000}) gegen vision-v2
    @44000 auf dem echten G1. Er entscheidet, ob der Echtbild-Verlust von N1.6 nach dem
    Co-Training tragbar ist; sonst ist N1.7 D die Alternative ohne diesen Verlust, das in der
    Simulation aber noch nicht greift. Voraussetzung ist eine Roboter-Seite nach
    \cref{sec:if-roboter}.
  \item \textbf{Griffrate heben.} Die vorregistrierte Regel sieht nach bestandenem Gate vor, den
    Weg aufzustocken, also mehr Episoden zu rendern; eine höhere Mischung ist nicht angezeigt,
    weil die Validierungs-\ac{mse} schon 3--9\,\% über Lauf~3 liegt (\cref{sec:erg-cotrain}). Alternativ Reinforcement Learning mit \ac{fpo}~\cite{fpo2025}, das bei flow-basierten \acp{vla} in Simulation bereits Verbesserungen gezeigt hat~\cite{pirl2025}, auf
    Modell C; der Trainer ist gebaut (nur N1.6), mit 3\,\% Griffen gibt es erstmals ein
    schwaches Startsignal (\doc{docs/weiterfuehrend/rl-anleitung.md}).
  \item \textbf{Gate-Schwelle nachschärfen:} relativ zum besten nicht co-trainierten Modell
    derselben Generation, Echtbild-Kriterium als feste Zahl.
  \item \textbf{\texttt{lifted\_clear} ohne Videobefund:} ein Kontaktkriterium, das einen
    gestützten Würfel ausschließt.
  \item \textbf{Ursachen offen:} 3-cm-Versatz der Würfellage (Lauf~63) und die Hänger im
    Closed Loop (umgangen mit \texttt{tools/sim\_cl\_watchdog.sh}).
\end{enumerate}

\section{Bekannte Defekte und Stolperstellen}
\label{sec:defekte}

\begin{itemize}
  \item Die Prüfbilder des visuellen Gates und die Layout-Datei
    \texttt{layout\_heldout\_dy3.json} liegen nur auf dem \acs{ikr}-Server. Auf einem neuen Server
    werden sie mit \texttt{server\_rl\_run.sh layout}, \doc{tools/layout_shift.py} und
    \texttt{server\_rl\_run.sh render} neu erzeugt (\doc{docs/quickstart.md}, §2).
  \item Das Sim-Image auf Docker Hub muss nach Änderungen an \texttt{Simulation/} neu gebaut
    werden (\texttt{Simulation/update\_sim\_image.sh --standalone}).
\end{itemize}

\section{Nicht weiterverfolgt}
\label{sec:nicht-weiter}
\label{sec:locomotion}
\label{sec:rl-weiterfuehrend}

Für die generative Augmentierung gibt es einen eigenen Container: Er variiert echte
Trainingsvideos mit Cosmos-Transfer2.5~\cite{cosmostransfer2025} (Kanten-ControlNet,
Beleuchtung, Hintergrund, Tischfarbe) und baut daraus einen Datensatz, den das Co-Training
direkt einliest (\texttt{USE\_COTRAIN}). Image-Build und erste Inferenzläufe gelangen Ende
August 2026; eine Vorabprüfung fängt Abweichungen der Framezahl ab. Ein vollständiger
augmentierter Datensatz und ein Trainingslauf damit fehlen, weil kein Docker-Server mit den
nötigen rund 65~GB GPU-Speicher verfügbar war.
Die Augmentierung variiert echte Bilder und zielt damit eher auf die Robustheit am Roboter als
auf den Gap zur Simulation (\doc{docs/augmentation/anleitung.md}). Der extern aufgenommene Sim-Datensatz
\texttt{Fichtl00/Cube\_Stacking\_synth} (16 Episoden, in der Simulation teleoperiert) passt
schematisch nicht (57 statt 28 Zustandsdimensionen, 50 statt 30~fps); ein Co-Training-Lauf mit
der umgewandelten Fassung liegt nicht vor
(\doc{docs/training/synth-datensatz.md}). Für Fortbewegung empfiehlt die Recherche, eine
vorhandene Lauf-Policy mit der Manipulations-Policy zu kombinieren
(\doc{docs/weiterfuehrend/lokomotion-recherche.md}).
